\documentclass[12pt,a4paper]{article}
\usepackage[utf8]{inputenc}
\usepackage[T1]{fontenc}
\usepackage{amsmath,amsfonts,amssymb}
\usepackage{graphicx}
\usepackage{cite}
\usepackage{url}
\usepackage{geometry}
\usepackage{fancyhdr}
\usepackage{setspace}
\usepackage{booktabs}
\usepackage{algorithm}
\usepackage{algorithmic}
\usepackage{xcolor}

% Page setup
\geometry{margin=1in}
\doublespacing
\pagestyle{fancy}
\fancyhf{}
\rhead{\thepage}
\lhead{Online Social-Compliance Verification for AVs}

% Helper: planned (NSFC) content not yet supported by results
\newcommand{\planned}[1]{{\color{blue!60!black}\textit{[PLANNED: #1]}}}
% Helper: figure placeholder box (keeps the file compilable without image files)
\newcommand{\figph}[2]{\fbox{\parbox[c][#1][c]{0.92\linewidth}{\centering\itshape #2}}}

% Title and author information
\title{Online runtime verification of socially compliant autonomous driving}
\author{
    Author Name\thanks{Corresponding author: author@institution.edu} \\
    \textit{Department, Institution} \\
    \textit{City, Country}
}
\date{\today}

\begin{document}

\maketitle

\begin{abstract}
\noindent
Autonomous vehicles are certified for collision safety, yet they operate within human
traffic where competent behaviour is governed by implicit, context-dependent social norms.
We recast social compliance as online runtime verification: given an interaction, we ask
whether a vehicle's behaviour falls within the range that competent humans exhibit in the
same situation. Using 38{,}228 multi-source human interactions, we first show that social
compliance is state-dependent: the priority-related interaction preference reverses sign
with collision risk ($+0.058$ to $-0.034$) and is robust across data sources and binning
choices, so a single sociality score is ill-posed. We then ask what an online verifier must
read. Counterintuitively, predicting risk and looking up an empirical envelope barely
narrows the reasonable interval, whereas the human reasonable range conditioned on the
driver's own early-window causal preference, calibrated by conformal prediction, narrows it
by $42\%$ at nominal $90\%$ coverage and, uniquely, transfers across a held-out data source
($0.902$). The
resulting verifier is sharp, calibrated and source-transferable, excludes risk from its
online inputs, and exposes a deviation signal to the planner, advancing auditable,
human-aligned autonomy.
\end{abstract}

% ======================================================================
\section{Introduction}
% ----------------------------------------------------------------------
Autonomous vehicles (AVs) must share the road with human drivers, cyclists and pedestrians
whose behaviour is shaped not only by physical constraints but also by implicit social
norms. A manoeuvre that is collision-free in isolation may still be perceived as aggressive,
hesitant or norm-violating when performed in the wrong interaction context, eroding the
predictability on which mixed human--machine traffic depends. Controlled experiments show
that machine agents embedded in human groups can suppress beneficial social norms such as
reciprocity, with collective consequences~\cite{shirado2023reciprocity}. Social competence
is therefore not a cosmetic add-on but part of what it means for autonomy to be safe in a
human world.

Two established families of methods address adjacent but distinct questions. Offline
behavioural evaluation asks whether an AV, on average, drove like a human, typically by
reporting a single aggregate sociality score over a dataset. Formal safety verification asks
whether a trajectory will collide or violate a hard constraint, and can guarantee legal
safety as an online layer around a planner~\cite{pek2020online}. Neither answers the
question a socially aware AV faces moment to moment: given the current interaction state, is
this behaviour within the range that competent humans exhibit? This question is local in
time, conditional on context, and graded rather than binary.

We formulate socially compliant driving as the runtime verification of online-inferred
interaction preferences. The monitored property is not a hand-written temporal-logic
specification but a data-driven empirical norm learned from large-scale human interaction,
expressed through the interaction preference value (IPV), a trajectory-derived measure of
how an agent balances its own progress against the group's. An online verifier estimates the
IPV, compares it with the human reasonable range for the current situation, and exposes the
result to the planner as a warning, soft cost, fallback trigger or audit record.

This framing differs from three existing paradigms. Unlike offline AV evaluation, it is
online, per-timestep and context-conditioned rather than aggregate. Unlike runtime
verification of formal specifications, the monitored property is an empirical normative
distribution estimated from human data. Unlike socially compliant planning, it does not
encode sociality into a planner's objective but provides a pluggable, planner-agnostic
monitor that any planner can be checked against.

Our contributions are fourfold. First, we formalise social compliance as a runtime
verification problem and instantiate it as an online verifier. Second, we show that social
compliance is state-dependent: the social meaning of right-of-way reverses with collision
risk, robustly across four independent data sources, so compliance cannot be a scalar.
Third, and counterintuitively, we show that the verifier should not estimate risk: the
reasonable IPV interval is set by the driver's own early-window causal preference, and a
self-anchored, conformally calibrated verifier is sharp, attains nominal coverage and
transfers across a held-out source, whereas risk-based envelope lookup does not. Fourth,
\planned{external validation on an independent real-vehicle challenge dataset, and a
demonstration that social-compliance deviation carries information about outcomes that
collision-safety checks miss}. The present manuscript reports the first three contributions
in full and scopes the fourth as planned validation.

% ======================================================================
\section{Results}
% ----------------------------------------------------------------------

\subsection{Social compliance is a state, not a score}
We analysed 38{,}228 interaction cases from four naturalistic sources within the InterHub
corpus (Waymo 23{,}218; nuPlan 7{,}499; Lyft 5{,}105; Argoverse-2 2{,}406; 3{,}695{,}981
trajectory frames). For each interacting agent we computed the IPV, parameterised as an
angle $\theta\in[-\pi/2,\pi/2]$, where $\theta>0$ denotes cooperative behaviour.

The social meaning of right-of-way is gated by collision risk (Fig.~\ref{fig:state-dep}).
The difference in IPV between the priority and non-priority agent reverses sign as
post-encroachment time (PET) increases: under high risk (PET\,$\le$\,1.0\,s) the priority
agent is markedly more prosocial ($+0.058$, 95\%\,CI $[+0.050,+0.067]$, $n=9{,}799$); at
intermediate risk the gap vanishes ($+0.001$, $n=22{,}833$); and under low risk
(PET\,$>$\,2.0\,s) it inverts, the priority agent becoming relatively more competitive
($-0.034$, $[-0.045,-0.023]$, $n=4{,}863$). Right-of-way is thus a risk-modulated
negotiation variable rather than a fixed disposition. A coarse geometric prior is stable
across sources, and a hierarchical partial-pooling estimator recovers the per-state norms
with leave-one-dataset-out median absolute error $0.142$. The direction of the risk-gated
effect is preserved under dropping Waymo, under three-, five- and quantile-binning of PET,
and under geometry coarsening (Extended Data). Social compliance therefore cannot be reduced
to a single average; it is a contextual property.

\subsection{The reasonable interval is conditioned on the driver's own behaviour, not on risk}
Establishing that the norm is contextual does not by itself tell us what an online verifier
should read. The natural approach is to predict the situation's risk and look up the
corresponding empirical envelope. This adds little. On held-out human interactions, even
\emph{oracle} knowledge of PET narrows the reasonable IPV interval by only about $3\%$
relative to a global interval (mean width $0.833$ versus $0.857$ at matched $\approx$$90\%$
coverage; Fig.~\ref{fig:reframe}). Risk bins are not the principal source of information
about how tightly an individual's behaviour is bounded.

The driver's own early behaviour is. Conditioning the interval on the agent's early-window
causal rolling-IPV, reconstructed from a trailing trajectory prefix and a static
map-lane-centreline reference, narrows it to $0.485$ (a $42\%$ reduction) while holding
coverage near nominal ($0.901$); kinematics without the self-anchor reach $0.627$. The
dominant uncertainty in any single interaction is between-driver, and a driver's preference
is temporally self-consistent, so the most informative covariate is the agent's own recent
IPV rather than the situation's risk. Crucially, the reasonable range itself remains the
\emph{human population's conditional distribution} of full-interaction IPV given the early
signal and context: the early-window IPV locates an agent within the human range, it is not a
licence for an agent to define its own norm, and the range is calibrated to cover human
outcomes by split-conformal prediction (necessary because raw quantile models under-cover,
around $0.86$). Whether conditioning on an agent's own behaviour could mask a persistently
non-compliant agent is an explicit validity question we treat in the limitations and planned
validation.

\subsection{A calibrated, source-transferable online verifier}
The decisive test is transfer across data sources. On a locked balanced slice, holding out
the entire Waymo source for evaluation, the self-anchored interval is the only method that
reaches nominal coverage ($0.902$, Wilson 95\%\,CI $[0.894,0.910]$; width $0.628$); a global
interval ($0.868$), oracle-PET envelope ($0.860$) and self-anchor-free kinematics ($0.857$)
all under-cover (Fig.~\ref{fig:transfer}). Cross-source generalisation is therefore carried
by the driver's own early-window preference, not by risk binning or generic kinematics.

This early-window signal is causal rather than a leakage artifact. Reconstructed against an
observed-trajectory-prefix reference, the causal rolling-IPV correlates only weakly with the
offline IPV (corr $0.281$); reconstructed against a static map-lane-centreline reference,
which is route-conditioned and available online, the correlation rises to $0.993$ (MAE
$0.027$; Extended Data). The self-anchor signal thus derives from the trajectory prefix and
static map geometry, under the deployment assumption that the lane or route is known at
decision time.

Replacing the empirical-envelope lookup with the self-anchored interval in an otherwise
identical verifier, behind the same deviation-score interface, improves the verifier
(Fig.~\ref{fig:ab}). In-distribution, coverage is unchanged while the interval narrows
($0.485$ versus $0.833$) and the false-flag rate does not rise ($0.052$ versus $0.051$).
Under the held-out source, the self-anchored verifier improves coverage ($0.823$ versus
$0.786$), narrows the interval ($0.488$ versus $0.678$) and lowers the false-flag rate
($0.114$ versus $0.142$). The held-out coverage of the integrated verifier remains below
nominal, so the supported deployment uses the signal as a soft cost, warning or monitor, and
reserves hard constraints for after target-domain recalibration.

\subsection{Verifier output as a planner-facing channel}
The calibrated interval yields a signed deviation score that maps to planner-facing
outputs: a monitor record, a one-sided soft cost active only where the norm expects
prosociality, a warning, and a fallback candidate (Fig.~\ref{fig:planner}). As a second,
joint check, the verifier also evaluates whether the pair-sum and pair-difference of the two
agents' IPVs fall within their human reasonable ranges, capturing configurations such as
mutual competition that are individually plausible but jointly anomalous; this paired check
is a gate rather than a stronger primary predictor. Injecting a counterfactual
more-competitive violation into $4{,}800$ case-agents raises the fallback-trigger rate from
$3.0\%$ to $12.9\%$ and lifts the median pre-onset cost by $+0.252$, while leaving compliant
streams largely unpenalised (mean social cost $0.134$, false-trigger rate $3.0\%$). This is
an interface demonstration of actionability, not a closed-loop planner result.

\subsection{External validation on an independent real-vehicle challenge}
\planned{We will run the verifier on the award-winning algorithms of a national real-vehicle
autonomous-driving challenge, whose official rankings, per-scenario scores and safety-event
records provide an independent outcome ground truth. The analysis will test criterion
validity (whether higher-ranked algorithms sit closer to the human reasonable interval), the
consequence chain (whether higher deviation predicts more safety events and lower scores;
Fig.~\ref{fig:nsfc-validation}), and discriminant value: whether social-compliance deviation
retains predictive power for independent scores and events after controlling for
conventional safety and rule-compliance metrics, including cases that pass safety
verification yet are flagged socially (Fig.~\ref{fig:discriminant}). This establishes the
verifier as complementary to, not a relabelling of, collision-safety verification.}

% ======================================================================
\section{Discussion}
% ----------------------------------------------------------------------
We have reframed AV social behaviour from an offline evaluation metric into an online
runtime-verification property, and identified what such a verifier should read. Social
compliance is contextual: the human norm's location shifts with collision risk and geometry.
Yet for verifying an individual online, the informative signal is not the situation's risk
but the driver's own early-window preference. Anchoring on that signal and calibrating it
with conformal prediction yields intervals that are simultaneously sharp, nominally covered
and transferable across a held-out source, the property that risk-based and global
references both fail to achieve. Because the verifier deliberately excludes risk and PET from
its online inputs, it cannot be a relabelling of collision-safety verification; it measures a
distinct, complementary property.

Our claims are bounded. The verifier is an empirical, probabilistic monitor, not a formal
proof, and complements rather than replaces reachability-based safety
layers~\cite{pek2020online}. The self-anchored path is deployable on roughly $74\%$ of cases,
where a lane or route reference is available; the remainder fall back to a self-anchor-free
kinematic interval. The strong cross-source result holds on a balanced, lane-referenced
locked slice and is not an unconditional guarantee for arbitrary maps or routes, and the
integrated verifier still under-covers under source shift, so hard-constraint deployment
requires target-domain recalibration. A rival explanation is that the IPV re-encodes
kinematics such as relative speed or passing order rather than a genuine social preference;
the lane-reference causal reconstruction and self-anchor-free controls argue against a pure
artifact, but kinematic-only and IPV-ablated comparisons are needed to exclude it. A central
and related concern is norm-laundering: conditioning the reasonable range on an agent's own
early behaviour could in principle make a persistently non-compliant agent appear normal. The
range is the human population's conditional distribution rather than the agent's own, but
confirming that it still flags self-consistent, human-implausible behaviour, that it
generalises to unseen drivers, and that its verdicts track independent outcomes requires
dedicated validation (cross-individual coverage, a consistent-deviator stress test, and
adjudication of self-anchored versus situational verdicts against downstream outcomes), which
we set out as planned work. Beyond
driving, the framework offers a template for monitoring whether autonomous agents behave
within human normative ranges online, under uncertainty and distribution shift, rather than
optimising a single average behavioural metric.

% ======================================================================
\section{Methods}
% ----------------------------------------------------------------------

\subsection{Data and interaction extraction}
We used 38{,}228 interaction cases from four sources within InterHub (Waymo, nuPlan, Lyft,
Argoverse-2), comprising 3{,}695{,}981 trajectory frames. Each case identifies a dyadic
interaction with an annotated priority (right-of-way) agent; the priority label denotes the
right-of-way agent identity, not a numerical rank or the realised passing order. The
interval-estimation results use a balanced locked slice of 5{,}000 cases (10{,}000
agent-rows) on which a lane or route reference is available.

\subsection{Interaction preference value and causal self-anchor}
The IPV quantifies an agent's balance between self-interest and group reward, parameterised
as $\theta\in[-\pi/2,\pi/2]$. For online use we estimate a \emph{dynamic} IPV from a trailing
trajectory prefix (history window of order $1$--$2$\,s) using a real-time estimator with a
static map-lane-centreline reference, yielding a strictly causal self-anchor signal. The
choice of reference is decisive for causality: against an observed-prefix reference the
causal estimate correlates only weakly with the offline IPV (corr $0.281$), whereas against
the static map-lane reference it agrees closely (corr $0.993$, MAE $0.027$), confirming that
the signal derives from the prefix and static geometry rather than future information.
\emph{[TO SPECIFY: exact history-window length; estimator configuration.]}

\subsection{Reasonable-interval estimation and calibration}
We estimate the human reasonable interval of the (full-window) IPV by conditional quantile
regression on the self-anchored signal together with geometry, right-of-way role, AV/HV
identity, kinematics and the counterpart's lagged IPV; PET and other risk variables are
\emph{excluded} from the online feature path. Model quantiles are turned into a calibrated
interval at target coverage by split-conformal prediction. When source health is degraded or
the data source is unknown, the system switches to a more conservative source-guard
calibration radius; when a lane reference is unavailable, it falls back to a self-anchor-free
causal-kinematics interval. The signed nonconformity of the current IPV relative to the
calibrated interval is the deviation score exposed to the planner.

\subsection{Paired joint reasonableness}
As a second-layer check the verifier evaluates the pair-sum and pair-difference of the two
agents' IPVs against their human reasonable ranges. Empirically the paired signal is a joint
gate rather than a stronger primary predictor (a paired-only interval improves coverage by
under half a percentage point while widening the interval), so it is used to flag jointly
anomalous configurations, not to replace the self-anchored marginal check.

\subsection{Leakage contract}
Dynamic self-anchored IPV, current kinematics and static map-lane references are
strict-online and admissible as verifier inputs; observed PET, realised passing order,
full-window IPV and post-hoc phase labels are offline-only, used solely for calibration and
evaluation, never as deployed inputs.

\subsection{Evaluation protocols}
We report coverage, mean interval width and the Winkler interval score, with Wilson
confidence intervals on coverage. Cross-source robustness uses leave-one-dataset-out
evaluation, holding out the entire Waymo source. The verifier is evaluated head-to-head
against the empirical-envelope lookup behind a shared deviation-score interface. Planned
ablations include kinematic-only and IPV-removed variants that test whether the social signal
is separable from raw kinematics.

\subsection{Planned external validation protocol}
\planned{The verifier will be applied to the trajectories of award-winning algorithms from a
national real-vehicle challenge. We will align scenarios, compute per-algorithm and
per-scenario deviation, and relate deviation to official rankings, per-scenario scores and
safety events via correlation and incremental regression that controls for conventional
safety and rule-compliance metrics; where human reference trajectories are available, AV
algorithms and humans will be placed within the same reasonable interval.}

% ---------------------------------------------------------------------
\begin{algorithm}[t]
\caption{Offline calibration of the self-anchored reasonable interval}
\label{alg:calib}
\begin{algorithmic}[1]
\REQUIRE human interactions $\mathcal{D}$; history window; target coverage $1-\alpha$
\ENSURE quantile model $\hat{Q}$, conformal radius $\kappa$, source-guard radii
\STATE split $\mathcal{D}$ into estimation set $\mathcal{D}_e$ and calibration set $\mathcal{D}_k$
\FOR{each interaction in $\mathcal{D}_e$}
    \STATE $\theta^{\mathrm{self}}\leftarrow$ causal rolling-IPV from prefix + map-lane reference
    \STATE features $\leftarrow$ ($\theta^{\mathrm{self}}$, geometry, role, AV/HV, kinematics, counterpart lagged IPV) \COMMENT{no PET}
    \STATE target $\leftarrow$ full-window IPV
\ENDFOR
\STATE fit conditional quantile model $\hat{Q}_{[q_{05},q_{50},q_{95}]}$ on (features, target)
\STATE on $\mathcal{D}_k$: $\kappa\leftarrow$ split-conformal radius for coverage $1-\alpha$
\STATE per source: source-guard radius from source-conditioned calibration
\RETURN $\hat{Q},\,\kappa,\,\text{source-guard}$
\end{algorithmic}
\end{algorithm}

% ---------------------------------------------------------------------
\begin{algorithm}[t]
\caption{Online social-compliance verification step at time $t$}
\label{alg:online}
\begin{algorithmic}[1]
\REQUIRE prefix window $W_t$; map-lane reference; calibrated $\hat{Q},\kappa$; source health
\ENSURE reasonable interval, deviation score, joint flag, planner signal
\IF{lane/route reference available}
    \STATE $\theta^{\mathrm{self}}_i\leftarrow$ causal rolling-IPV from $W_t$ + map-lane \COMMENT{strict-online}
\ELSE
    \STATE fall back to self-anchor-free causal-kinematics features \COMMENT{$\approx$26\% of cases}
\ENDIF
\STATE features $\leftarrow$ ($\theta^{\mathrm{self}}_i$, geometry, role, AV/HV, kinematics, $\theta_j$ lagged) \COMMENT{PET excluded}
\STATE mode $\leftarrow$ source-guard \textbf{if} source unknown/drift \textbf{else} CQR
\STATE $[q_{05},q_{50},q_{95}]\leftarrow \hat{Q}(\text{features})$ widened by $\kappa$(mode) \COMMENT{conformal interval}
\STATE $\delta\leftarrow$ signed nonconformity of current $\theta_i$ vs interval
\STATE joint $\leftarrow$ (pair-sum and pair-diff of $\theta_i,\theta_j$ within human bands)
\IF{\textbf{not} \textsc{QualityGates}(support, provenance, source health)}
    \RETURN \textsc{Monitor}($\delta$, joint)
\ENDIF
\STATE signal $\leftarrow$ \textsc{Warning}+soft cost \textbf{if} $\delta$ exceeds threshold \textbf{else} \textsc{Monitor}
\RETURN (interval, $\delta$, joint, signal)
\end{algorithmic}
\end{algorithm}

\noindent\textbf{Complexity.} Each online step is $O(1)$ given the fitted quantile model and
conformal radius (one estimator call, one interval evaluation, a scalar deviation and a
paired check), so the verifier runs within a planning cycle and is auditable.

% ======================================================================
% Figures (placeholders kept compilable; replace \figph with \includegraphics)
% ======================================================================
\begin{figure}[t]\centering
\figph{3.0cm}{Figure 1 placeholder. Online verifier architecture (self-anchored path).}
\caption{\textbf{Online social-compliance verification framework.} A real-time estimator
reconstructs the driver's causal early-window IPV from a trajectory prefix and a static
map-lane reference; geometry, role, AV/HV identity and kinematics condition a
conformally-calibrated reasonable interval, with PET excluded from the online path. The
deviation score, a paired joint check, a source-guard calibration switch and a lane-fallback
branch feed the planner interface.}
\label{fig:architecture}
\end{figure}

\begin{figure}[t]\centering
\figph{3.0cm}{Figure 2 placeholder. State dependence of the population norm.}
\caption{\textbf{Social compliance is state-dependent.} The priority-minus-non-priority IPV
gap reverses with PET risk ($+0.058\!\to\!+0.001\!\to\!-0.034$) with 95\% confidence
intervals and per-source forest plots, and a coarse geometry prior is prosocial across all
four sources. A single sociality score is ill-posed.}
\label{fig:state-dep}
\end{figure}

\begin{figure}[t]\centering
\figph{3.0cm}{Figure 3 placeholder. Source: fig1\_reframe.}
\caption{\textbf{Risk barely narrows the reasonable interval; the driver's own preference
does.} On held-out interactions, oracle PET narrows the interval by only $\approx$$3\%$
relative to a global interval ($0.833$ versus $0.857$), whereas causal rolling-IPV narrows it
to $0.485$ ($-42\%$), with coverage near nominal in all cases.}
\label{fig:reframe}
\end{figure}

\begin{figure}[t]\centering
\figph{3.0cm}{Figure 4 placeholder. Source: fig2\_cross\_dataset.}
\caption{\textbf{Only the self-anchored interval transfers across sources.} With Waymo held
out (locked balanced slice), causal rolling-IPV is the only method reaching nominal $0.90$
coverage ($0.902$, Wilson 95\%\,CI $[0.894,0.910]$); global, oracle-PET and self-anchor-free
baselines under-cover.}
\label{fig:transfer}
\end{figure}

\begin{figure}[t]\centering
\figph{3.0cm}{Figure 5 placeholder. Source: fig4\_verifier\_ab.}
\caption{\textbf{Verifier A/B: self-anchored interval versus empirical-envelope lookup.}
Behind the same deviation interface, the self-anchored verifier is markedly narrower in and
out of distribution and does not increase false flags; under the held-out source it improves
coverage and lowers the false-flag rate ($0.142\to0.114$).}
\label{fig:ab}
\end{figure}

\begin{figure}[t]\centering
\figph{3.0cm}{Figure 6 placeholder. Planner-facing interface demonstration.}
\caption{\textbf{Planner-facing soft-cost demonstration.} Counterfactual violation injection
raises the fallback-trigger rate from $3.0\%$ to $12.9\%$ and the median pre-onset cost by
$+0.252$, while leaving compliant streams largely unpenalised. An interface demonstration,
not a closed-loop or safety guarantee.}
\label{fig:planner}
\end{figure}

\begin{figure}[t]\centering
\figph{2.8cm}{Figure 7 placeholder \planned{NSFC external validation}}
\caption{\planned{\textbf{External validation on an independent real-vehicle challenge.}
Per-algorithm social-compliance deviation versus official ranking and per-scenario scores;
deviation versus safety-event rate.}}
\label{fig:nsfc-validation}
\end{figure}

\begin{figure}[t]\centering
\figph{2.8cm}{Figure 8 placeholder \planned{NSFC discriminant value}}
\caption{\planned{\textbf{Discriminant value beyond safety verification.} Cases that pass
safety verification yet are flagged by the social verifier, and incremental prediction of
scores and events after controlling for conventional safety metrics.}}
\label{fig:discriminant}
\end{figure}

% ======================================================================
\section*{Data Availability}
InterHub-derived interaction cases and source datasets (Waymo, nuPlan, Lyft, Argoverse-2)
are governed by their respective licences; processed derivatives will be made available upon
publication. \planned{Real-vehicle challenge data availability to be stated per the
challenge organiser's terms.}

\section*{Code Availability}
Implementation, the real-time IPV estimator, leakage-contract unit tests and
figure-generation code will be released upon publication.

\section*{References}
\bibliographystyle{unsrt}
\bibliography{bibliography/biblio}

\section*{Acknowledgements}
% (Funding sources, institutional support, data providers)

\section*{Author Contributions}
% (Individual contributions following Nature journal format)

\section*{Competing Interests}
The authors declare no competing interests.

\section*{Additional Information}
\subsection*{Supplementary Information}
\subsection*{Supplementary Tables}
\subsection*{Supplementary Figures}

\end{document}
